\chapter{Software Requirements Specification}

A Software Requirements Specification (SRS) records what a software system must do and the qualities it must have. Functional requirements describe the behaviour the system provides, and non-functional requirements describe properties such as performance, reliability, usability and security \cite{sommerville2016se,iso29148}. A clear SRS gives all stakeholders a shared reference: developers know what to build, testers know what to verify, and supervisors can judge progress against agreed expectations. It also limits scope creep, because new ideas can be checked against what was actually agreed.

For a project developed in increments, the SRS also shows which requirements have already been met. This chapter introduces the document's scope and audience, lists the functional requirements of the implemented modules and the planned system with their current status, defines the non-functional requirements, describes the techniques used to gather the requirements, and presents the project time frame.
